% =============================================================================
% 048 / 068
% Status: raw
% =============================================================================
% Notes:
%
% Source question:
% 漢字為什麼會導致理性思辨崩毀呢？當中的運作機理和前因後果是什麼？
% =============================================================================

\chapter[漢字為什麼會導致理性思辨崩毀呢？當中的運作機理和前因後果是什麼？]{漢字為什麼會導致理性思辨崩毀呢？當中的運作機理和前因後果是什麼？}
\qaanswer{
漢字不會導致邏輯思辨出現障礙，但文字用純漢字表述會導致邏輯思辨的混亂。理由如下：

1. 漢字基本都有實意，當漢字作為虛詞成分出現時，它的基本實意會干擾解讀者理解它作為語法標記時的作用和意思。全漢字文字讓虛詞、實詞的區分和識別變得困難；

2.虛詞、實詞難以區別和辨識就會導致語法建構出現問題，難以形成自洽的語法體系，甚至讓很多漢語母語者以為漢語沒有語法；

3.漢字的延展性導致了大量的望文生義，特別是成語、習慣語的字面意思模糊，非常容易導致不同時代背景的解讀者朝著相反的方向理解，著名的例子有「七月流火」、「空穴來風」、「不孝有三、無後為大」。

以上這幾個問題相互交織疊加讓漢語的語言思維更是變得混亂不堪，直接導致漢語母語者多數會在邏輯思辨層面上出現障礙。
}

